% Part: second-order-logic
% Chapter: sol-and-sets
% Section: comparing-sets

\documentclass[../../../include/open-logic-section]{subfiles}

\begin{document}

\olfileid{sol}{set}{cmp}

\olsection{Mbandhingake Himpunan}

\begin{prop}
!!{formula} $\lforall[x][(X(x) \lif Y(x))]$ netepake relasi
subhimpunan, yaiku
$\Sat{M}{\lforall[x][(X(x) \lif Y(x))]}[s]$
yen lan mung yen $s(X) \subseteq s(Y)$.
\end{prop}

\begin{prop}
!!{formula} $\lforall[x][(X(x) \liff Y(x))]$ netepake relasi
identitas ing himpunan, yaiku
$\Sat{M}{\lforall[x][(X(x) \liff Y(x))]}[s]$
yen lan mung yen $s(X) = s(Y)$.
\end{prop}

\begin{prop}
!!{formula} $\lexists[x][X(x)]$ netepake sipat ora kosong,
yaiku $\Sat{M}{\lexists[x][X(x)]}[s]$ yen lan mung yen
$s(X) \neq \emptyset$.
\end{prop}

Himpunan~$X$ ora luwih gedhe tinimbang himpunan~$Y$,
$\cardle{X}{Y}$, yen lan mung yen ana fungsi
!!a{injective} $f\colon X \to Y$. Amarga kita bisa
nyatakake yen fungsi iku injektif, lan uga yen nilai-nilaine
kanggo argumen ing~$X$ kalebu ing~$Y$, kita uga bisa
netepake relasi ora luwih gedhe tinimbang ing antarane
subhimpunan domain.

\begin{prop}
Rumus
\[
\lexists[u][(\lforall[x][(X(x) \lif Y(u(x)))] \land \lforall[x][\lforall[y][(\eq[u(x)][u(y)] \lif
      \eq[x][y])]])]
\]
netepake relasi ora luwih gedhe tinimbang.
\end{prop}

Rong himpunan padha ukurane, utawa ``padha cacahe,''
$\cardeq{X}{Y}$, yen lan mung yen ana fungsi
!!a{bijective}~$f\colon X \to Y$.

\begin{prop}
Rumus
\begin{multline*}
  \lexists[u][(\lforall[x][(X(x) \lif Y(u(x)))] \land {}\\
    \lforall[x][\lforall[y][(\eq[u(x)][u(y)] \lif \eq[x][y])]]
  \land {} \\\lforall[y][(Y(y) \lif \lexists[x][(X(x)
      \land \eq[y][u(x)])])])]
\end{multline*}
netepake relasi padha cacahe.
\end{prop}

Kita bakal nyingkat !!{formula}s iki, kanthi urut,
dadi $X \subseteq Y$, $X = Y$, $X \neq \emptyset$,
$\cardle{X}{Y}$, lan $\cardeq{X}{Y}$.
(Iki bisa rada mbingungake, amarga kita nggunakake notasi
kang padha nalika ngomong kanthi ora formal ngenani
himpunan $X$ lan $Y$---nanging ing kene notasi iku
singkatan kanggo !!{formula}s logika orde kapindho
kang ngemot variabel relasi siji-papan $X$ lan~$Y$.)

\begin{prop}
!!{sentence} $\lforall[X][\lforall[Y][((\cardle{X}{Y} \land
    \cardle{Y}{X}) \lif \cardeq{X}{Y})]]$ iku valid.
\end{prop}

\begin{proof}
!!{sentence} iku dicukupi ing !!a{structure}~$\Struct{M}$
yen, kanggo saben subhimpunan $X \subseteq \Domain{M}$
lan $Y \subseteq \Domain{M}$, saka
$\cardle{X}{Y}$ lan $\cardle{Y}{X}$ nyusul
$\cardeq{X}{Y}$. Nanging iki bener kanggo
\emph{sembarang} himpunan $X$ lan $Y$---yaiku
Teorema Schr\"oder-Bernstein.
\end{proof}


\end{document}
